% ============================================================
%  COMANDOS PROPIOS Y CONVENCIONES DE ESCRITURA
% ============================================================
%
%  MATEMÁTICA
%   - En línea: $...$          (nunca \( ... \))
%   - En display: \[ ... \] para UNA línea; align* para varias, partiendo
%     en "&=" para que ninguna línea supere el ancho de la columna (~3.5in).
%     Nunca $$ ... $$ (sintaxis de TeX plano: amsmath no controla su espaciado).
%     Solo usar display cuando la fórmula no cabe en línea: cada display
%     cuesta ~3 líneas de columna.
%   - Complejidad: SIEMPRE \bigO{...} y el argumento en notación matemática:
%       \bigO{n}   \bigO{n \log n}   \bigO{V + E}   \bigO{V \cdot E}
%       \bigO{n^2}   \bigO{\pow{2}{k}}   \bigO{M \alpha(N)}
%     Sin frases dentro de \bigO (\text{...} no se parte de línea y desborda la
%     columna): usar una letra y explicarla después:
%       \bigO{\pow{2}{k} \cdot C}, con $C$ = costo de evaluar una combinación.
%     Solo palabras cortas: \bigO{\log(\text{rango})}
%   - Potencias de 10:  \ten{5}  -> 10^5      Otras potencias: \pow{2}{k} -> 2^k
%   - Aproximado:       \aprox{}0.1s -> ~0.1s
%   - Piso/techo:       \floor{t/2}  \ceil{n/2}
%   - Los símbolos Unicode ≤ ≥ → − × · ² Σ se pueden usar en texto corrido
%     (ver packages.tex). Dentro de $...$ o \bigO{} usar \le \ge \to - \times \cdot.
%
%  TEXTO
%   - Etiquetas de cada entrada: \graybold{Complejidad:} etc.
%   - Identificadores / código en línea: \texttt{...} o \code{...}
%
% ------------------------------------------------------------
%  Matemática
% ------------------------------------------------------------
\newcommand{\bigO}[1]{\ensuremath{\mathcal{O}(#1)}}      % \bigO{n \log n}
\newcommand{\pow}[2]{\ensuremath{{#1}^{#2}}}            % \pow{2}{k}  -> 2^k
\newcommand{\ten}[1]{\ensuremath{10^{#1}}}              % \ten{5}     -> 10^5
\newcommand{\aprox}{\ensuremath{{\sim}}}                % \aprox{}2s  -> ~2s  (las llaves evitan el espacio de relación)
\newcommand{\set}[1]{\ensuremath{\{#1\}}}               % \set{1,2,3}
\newcommand{\abs}[1]{\ensuremath{\lvert #1 \rvert}}     % \abs{x}
\newcommand{\floor}[1]{\ensuremath{\lfloor #1 \rfloor}} % \floor{t/2}
\newcommand{\ceil}[1]{\ensuremath{\lceil #1 \rceil}}    % \ceil{t/2}

% ------------------------------------------------------------
%  Texto
% ------------------------------------------------------------
\newcommand{\graybold}[1]{\textcolor{gray}{\textbf{#1}}} % etiquetas de sección
\newcommand{\code}[1]{\texttt{#1}}                       % código en línea
\newcommand{\regex}[1]{\texttt{#1}}                      % expresión regular en línea

% ------------------------------------------------------------
%  Títulos (titlesec) — color y numeración arábiga plana
% ------------------------------------------------------------
\titleformat{\section}
  {\Large\bfseries\color{TitleBlue}} % formato
  {\thesection}                      % número
  {0.5em}                            % separación
  {}                                 % antes del título
\titleformat{\subsection}
  {\large\bfseries\color{TitleBlue}}
  {\thesubsection}
  {0.5em}
  {}
\titleformat{\subsubsection}
  {\normalsize\bfseries\color{TitleBlue}}
  {\thesubsubsection}
  {0.5em}
  {}
% IEEEtran numera con romanos (I, II-A, ...); se fuerza 1, 1.1, 1.1.1
\renewcommand{\thesection}{\arabic{section}}
\renewcommand{\thesubsection}{\thesection.\arabic{subsection}}
\renewcommand{\thesubsubsection}{\thesubsection.\arabic{subsubsection}}
